\section{Exact finite selection certificates}

The real-valued greedy estimate is convenient for asymptotics. For
specified finite dimensions, integrality gives a more precise computable
guarantee. This section describes a certificate using rational arithmetic
only. It computes guaranteed ranks, not actual optimizing translates of an
unspecified map.

Let $u,v$ be rational and retain the notation of
Section~\ref{sec:sharp-spanning}. Put
\[
 R=\left\lceil\frac{(1-u)^2}{1+v}d\right\rceil.
\]
At a current span of dimension $k$, the square-root estimate and
rank-nullity guarantee a gain of at least
\[
 g(k)=\max\left\{0,\ s-\left\lfloor
 su+s\sqrt{(1+v)k/d}\right\rfloor\right\}.
\]
In particular $g(k)\ge1$ when $k<R$. The actual rank can exceed a
previous lower certificate, so blindly iterating $k\mapsto k+g(k)$ on a
lower bound would require a monotonicity argument. We avoid that issue by
taking an envelope. Define
\[
 T(k)=\min\{R,k+g(k)\},\qquad
 H(q)=\min_{q\le k\le R}T(k),\qquad
 q_0=0,\quad q_{j+1}=H(q_j).
\]
If the actual rank is at least $R$, the target is already reached.
Otherwise the actual integer rank lies between $q_j$ and $R-1$, so the
next actual rank is at least $H(q_j)$. Thus $q_j$ is a rigorous lower
certificate after at most $j$ selected translates. Furthermore
$q_{j+1}>q_j$ whenever $q_j<R$, hence the certificate reaches $R$ in at
most $R$ steps. The table of $H$ is obtained from one backward pass over
the values of $T$.

\subsection{Eliminating all square-root comparisons}

To decide whether an integer $m$ is at most
$su+s\sqrt{(1+v)k/d}$, set $z=m-su$. The answer is yes if $z\le0$;
if $z>0$, it is yes exactly when
\[
 z^2\le s^2(1+v)k/d.
\]
Every quantity is rational. Binary search over $0\le m\le s$ suffices
because values above $s$ all give gain zero. This constructs the table in
$O(R\log(s+1))$ exact comparisons and $O(R)$ storage. There is no
floating-point eigenvalue calculation or numerical linear program in the
acceptance test.

If the source is restricted to a subspace of codimension $r_0$, the same
$j$ selected translates retain rank at least $q_j-jr_0$, by the restriction
lemma. This keeps the dimension loss from source conditions explicit.

\subsection{An illustrative comparison}

For $d=1000$, $s=20$, and $u=v=1/100$, the sharp universal ceiling is
$R=971$. The exact integer certificate reaches rank $840$ after $148$
translates and rank $971$ after $261$. The exponential greedy bound in
the source, with its $16\varepsilon$ cutoff and rate $2s/(5d)$, certifies
rank $840$ after $230$ translates. These figures compare sufficient
guarantees; they do not assert a lower bound on the optimal number of
translates for a particular family.

The supplied checker also instantiates the source errors
$u=5/b^4$, $v=9/b^4$, verifies rational equality examples, and checks the
constants used in the $20/b^4$ theorem. Its output records the exact inputs,
rank trajectories, and certified translate budgets. The illustrative
dimensions are not claimed to occur as admissible automorphic data.

There is a second dimension constraint:
$\rank(PB)\ge\lceil(1-u)s\rceil$. Consequently $R$ is the ceiling of the
optimal \emph{universal coefficient times $d$}, not necessarily the
exact minimum possible total rank for every fixed triple $(d,s,u,v)$.
The cyclic examples establish the sharpness actually claimed.
